\documentclass[../main.tex]{subfiles}

\begin{document}
\section{Generalization to $|a| \neq 2$} 
\subsection{Notations}

For each $n,a$ let 
\[
  f_{n,a} = x^n + ax + 1
\]
and carry all the notations from the previous section.
Assume that $|a| \neq 0$.

Fix an algebraic closure $\overline{\Q}$ of $\Q$ and for each prime $p$, fix an algebraic closure $\overline{\Q_p}$ of $\Q_p$ and an embedding $\overline{\Q} \hookrightarrow \overline{\Q_p}$ extending the natural embedding $\Q \hookrightarrow \Q_p$.
Denote by $G$ the absolute Galois group $\Gal\left( \overline{\Q}/\Q \right)$, and for each prime $p$, denote by $I\left( p \right) \subseteq G$ the absolute inertia group at $p$, i.e., the embedding into $G$ of the inertia group $\Gal\left( \overline{\Q_p}/\Q_p^{\un} \right)$, where $\Q_p^{\un}$ is the maximal unramified extension of $\Q_p$ inside $\overline{\Q_p}$.

\subsection{Strategy}
We will first lay down a few results that will provide us with a strategy to show that for any pair $n,a$ that satisfies different conditions, the Galois group of any factor of $f_{n,a}$ is symmetric.
\begin{lemma}
  \label{lem:inertia_group_acts_as_transposition}
  For any $p \not | a$, the inertia group $I\left( p \right)$ acts on $\Hom_{\Q}\left( \Q\left[ x \right] / \left( f_{n,a} \right), \overline{\Q} \right)$ either trivially or as a single transposition. 
\end{lemma}
\editorial{I think I saw somewhere a proof that covers this case, but I can't find it now. So I will try to adapt Martin's proof}
\begin{proof}
  If $I\left( p \right)$ acts trivially, we are done.
  Assume that it does not.
  It suffices to show that $I\left( p \right)$ acts on the roots of $f_{n,a}$ as a single transposition.
  If two roots are in the same orbit under $I\left( p \right)$, then they have the same reduction modulo $p$.
  Therefore, it suffices to show that $f_{n,a}$ has at most one pair of roots that are congruent modulo $p$, i.e. if we denote by $\overline{f_{n,a}} \in \F_p\left[ x \right]$ the reduction of $f_{n,a}$ modulo $p$, then it suffices to show that $f_{n,a}$ has exactly one root in $\overline{\F_p}$ of multiplicity larger than $1$ and that its multiplicity is in fact $2$.

  Firstly note that by \cite[Theorem 2]{Swan62} 
  \[
    \disc\left( f_{n,a} \right) = 
    \left( -1 \right)^{\frac{n(n-1)}{2}} \left(  
      n^n + \left( 1 - n \right)^{n - 1} \left( n - 1 \right)^{n-1} a^n
    \right)
  \]
  and hence, since $p | \disc\left( f_{n,a} \right)$ and since $p \nmid |a|$, it follows that $p$ divides $n$ iff it divides $\left( n-1 \right)$ and thus it divides neither.
  
  Next, since 
  \begin{align*}
    & f_{n,a}' = n x^{n-1} + a \\
    & f_{n,a}'' = n \left( n - 1 \right) x^{n-2}
  \end{align*}
  and since $p \nmid n \left( n - 1 \right), a$ it follows that $f'_{n,a}, f''_{n,a}$ have no shared root in $\overline{\F_p}$ and hence $f_{n,a}$ has no triple roots in $\overline{\F_p}$.

  Finally, note that any double root of $f_{n,a}$ is also a solution of the linear equation
  \[
    g_{n,a} := 
    n f_{n,a} - x f_{n,a}' = 
    nx^n + a n x + n - n x^n - a x = 
    a \left( n - 1 \right) x + n
  \]
  and hence since $p$ does not divide $a \left( n - 1 \right)$ there is exactly one double root of $f_{n,a}$ in $\overline{\F_p}$.
\end{proof}
\begin{corollary}
  \label{cor:ramification_only_over_primes_dividing_a}
  Any number field $K$ that embeds into $\Q\left[ x \right] / \left( f_{n,a} \right)$ is ramified only over primes dividing $a$.
\end{corollary}
\begin{proof}
  It suffices to show that $K$ is unramified above any prime $p$ that does not divide $a$.
  Let $p$ be such a prime.
  It suffices to show that $I\left( p \right)$ acts trivially on $\Hom_{\Q}\left( K, \overline{\Q} \right)$.
  If $I\left( p \right)$ acts trivially on $\Hom_{\Q}\left( \Q\left[ x \right] / \left( f_{n,a} \right), \overline{\Q} \right)$, we are done.
  Assume otherwise.
  Choose an embedding $\iota: K \hookrightarrow \Q\left[ x \right] / \left( f_{n,a} \right)$.
  Then the precompositio map 
  \[
    \iota^\ast : \Hom\left( \Q\left[ x \right] / \left( f_{n,a} \right), \overline{\Q} \right) \to 
    \Hom_{\Q}\left( K, \overline{\Q} \right)
  \]
  is a surjective $G$-equivariant map.
  Thus, $I\left( p \right)$ acts trivially on $\Hom_{\Q}\left( K, \overline{\Q} \right)$ since any non-trivial action of $I\left( p \right)$ on $\Hom_{\Q}\left( K, \overline{\Q} \right)$ would lift to an action on $\Hom_{\Q}\left( \Q\left[ x \right] / \left( f_{n,a} \right), \overline{\Q} \right)$ that would have at least $4$ non-fixed points, contradicting the fact that by Lemma \ref{lem:inertia_group_acts_as_transposition}, $I\left( p \right)$ acts on $\Hom_{\Q}\left( \Q\left[ x \right] / \left( f_{n,a} \right), \overline{\Q} \right)$ as a single transposition.
\end{proof}

\begin{lemma}
  A transitive permutation group that is generated by transpositions is the full symmetric group.
\end{lemma}
\begin{proof}
  \editorial{Sketch inspired by a discussion with Itamar Israeli:}

  There is a seuqnece of transpositions that connects any two elements.
  Given a product of a pair of transpositions 
  \[
    \left( a \; b \right) \left( b \; c \right) = \left( a \; b \; c \right)
  \]
  it can be shoretned to 
  \[
    \left( b \; c \right) \left( a \; b \right) \left( b \; c \right) = \left( a \; c \right) 
  \]
\end{proof}
\begin{corollary}
  \label{cor:primitive_group_with_transposition_is_full_symmetric}
  A primitive permutation group that contains a transposition is the full symmetric group.
\end{corollary}
\begin{proof}
  \editorial{Slightly inspired by Borys's thesis:}

  Let $H$ be a primitive permutation group that contains a transposition $\tau$.
  Let $K$ be the normal closure of $\left< \tau \right>$ inside $H$.
  Then the orbits of $K$ form a system of blocks for $H$.
  Therefore, $K$ is itself transitive and it is generated by transpositions and hence by the previouse lemma $K = S_n$ and hence $H = S_n$ as required.
\end{proof}
\begin{corollary}
  If the galois group $G_L$ of a field $L$ of degree $n$ contains a transposition and $L$ has no proper intermediate extensions, then $G_L = S_n$.
\end{corollary}
\begin{proof}
  According to Corollary \ref{cor:primitive_group_with_transposition_is_full_symmetric} it suffices to show that $G_L$ is primitive.
  Assume that it is not.
  Then there is a strictly smaller non-trivial $G$ set $S$ and a surjective $G$-equivariant map
  \[
    \pi : \Hom_{\Q}\left( \Q\left( \alpha \right), \overline{\Q} \right) \to S
  \]
  In particular, $\pi$ is an epimorphism of $G$-sets, and hence by the famous equivalnece of categories between $G$-sets and Etale algebras over $\Q$, $\pi$ corresponds to a monomorphism of Etale algebras over $\Q$, i.e. a field embedding 
  \[
    \iota : K \hookrightarrow L
  \]
  and since $S$ is non-trivial and strictly smaller, $K$ is a proper subextension of $L/\Q$, contradicting the assumption that $L$ has no proper intermediate extensions.
\end{proof}
Our strategy would be to find conditions on $n,a$ that would guarantee that the Galois group over $\Q$ of any factor $f$ of $f_{n,a}$ contains a transposition and that $\Q\left[ x \right]/ \left( f \right)$ has no proper intermediate extensions.
For the second part, we will rely on Corollary \ref{cor:ramification_only_over_primes_dividing_a} and different bounds on root discriminants of number fields.

\subsection{$a = \pm 1$}


\begin{lemma}
  $f_{n,a}$ is irreducible over $\Q$.
\end{lemma}
\begin{proof}
  It suffices to show that 
  \[
    \widehat{f_{n,a}} := x^n f_{n,a}\left( 1/x \right) = x^n + ax^{n-1} + 1 
  \]
  is irreducible over $\Q$.
  For any $x \in S^1 \subseteq \C$, 
  \[
    \left| a x^{n-1} \right| = \left| a \right| > 2 = \left|x^n\right| + \left|1\right| \geq \left|x^n + 1\right|  
  \]
  and hence the winding number of $\widehat{f_{n,a}}\left( x \right)$ around $0$ as $x$ traverses the unit circle is equal to that of $- a x^{n-1}$ which is $n-1$.
  Therefore, $\widehat{f_{n,a}}$ has $n-1$ roots in the interior of the unit disk and hence it is irreducible over $\Q$.
\end{proof}


\begin{lemma}
  \label{lem:valuation_of_disc}
  Let $p$ be a prime dividing $a$.
  If 
  \[
    v_p\left( n \right) \neq v_p\left( a \right) \text{ or } \left( n / a \right)^n \not \equiv -1 \mod p
  \]
  then 
  \[
    v_p\left( \disc\left( f_{n,a} \right) \right) = n \min\left\{  
      v_p\left( n \right), v_p\left( a \right) 
    \right\}
  \] 
\end{lemma}
\begin{proof}
  Assume first that $v_p\left( n \right) \neq v_p\left( a \right)$.
  According to \cite[Theorem 2]{Swan62}
  \[
    \disc\left( f_{n,a} \right) = 
    \left( -1 \right)^{\frac{n(n-1)}{2}} \left(  
      n^n + \left( 1 - n \right)^{n - 1} a^n
    \right)
  \]
  Then either $v_p\left( n \right) = 0$ and hence $v_p\left( \disc\left( f_{n,a} \right) \right) = 0$ or $v_p\left( n \right) > 0$ and hence $v_p\left( n - 1 \right) = 0$ and thus $v_p\left( \left( 1 - n \right)^{n - 1} a^n \right) = n v_p\left( a \right)$ and hence 
  \[
    v_p\left( \disc\left( f_{n,a} \right) \right) = 
    \min\left\{ v_p\left( n^n \right), v_p\left( \left( 1 - n \right)^{n - 1} a^n \right) \right\} = 
    n \min{v_p\left( n \right), v_p\left( a \right) }
  \]
  
  Assume now that $v_p\left( n \right) = v_p\left( a \right)$.
  Then 
  \[
    \disc\left( f_{n,a} \right) = 
    \Res\left( f_{n,a}, f_{n,a}' \right) = 
    p^{n v_p\left( a \right)} \Res\left(  
      f_{n,a},
      p^{-v_p\left( a \right)} f_{n,a}'
    \right)
  \]
  Since $p^{-v_p\left( a \right)} f_{n,a}' \in \Z\left[ x \right]$ it follows that $v_p\left( \disc\left( f_{n,a} \right) \right) \geq n v_p\left( a \right)$ and equality holds iff $f_{n,a}$ and $p^{-v_p\left( a \right)} f_{n,a}'$ have no common root in $\overline{\F_p}$.
  Assume then that there is no equality, i.e. that 
  \begin{align*}
    & \overline{f_{n,a}} = x^n + ax + 1 = x^n + 1 \in \F_p\left[ x \right] \\
    & p^{-v_p\left( a \right)} \overline{f_{n,a}}' = n p^{-v_p\left( a \right)} x^{n-1} + a p^{-v_p\left( a \right)} \in \F_p\left[ x \right]
  \end{align*}
  have a common root.
  This root in particulat annahilates 
  \[
    g_{n,a} := 
    x \overline{f_{n,a}}' - n \overline{f_{n,a}} = 
    n x^n + a x - nx^n - n = 
    a x - n
  \]
  i.e. it is equal to $n / a$ and hence
  \[
    0 = \overline{f_{n,a}}\left( n / a \right) = \left( n / a \right)^n + 1
  \]
  and thus $\left( n / a \right)^n \equiv -1 \mod p$.
\end{proof}
Assume in the sequel that $|a| < 22$ (or $< 44$ assuming GRH).
According to \cite{OdlyzkoSurvey}, there is a constant $C\left( \left| a \right| \right)$ such that for any number field $K$ with degree $m > C\left( \left| a \right| \right)$,
\[
  \left| \disc\left( K \right)^{1/m} \right| > \left| a \right|
\]
Moreover, let $S\left( \left| a \right| \right)$ be the set of positive integers $n$ that satisfy the conditions of Lemma \ref{lem:valuation_of_disc}, i.e. for any prime $p\mid a$ either 
\[
  v_p\left( n \right) \neq v_p\left( a \right) \text{ or } \left( n / a \right)^n \not \equiv -1 \mod p
\]
\begin{lemma}
  \label{lem:there_is_a_prime_dividing_disc_but_not_a}
  Let $C\left( \left| a \right| \right) < n \in S\left( \left| a \right| \right)$.
  Then there is a prime $p$ s.t. $p\mid \disc\left( \Q\left( \alpha \right) \right)$ but $p \not \mid a$ and consequently $\Q\left( \alpha \right)$ ramifies above $p$. 
\end{lemma}
\editorial{The only example I found so far where every prime dividing the discriminant also divides $a$ is $n = 3, a= -3$ and I checked for any $n \leq 25$ and $\left|a\right| \leq 22$.}
\begin{proof}
  According to Lemma \ref{lem:valuation_of_disc} for any prime $p$ dividing $a$ 
  \[
    v_p\left( \disc\left( \Q\left( \alpha \right) \right) \right) \leq
    v_p\left( \disc\left( f_{n,a} \right) \right) = 
    n \min\left\{ v_p\left( n \right), v_p\left( a \right) \right\} \leq
    n v_p\left( a \right)
  \]
  and hence if any prime dividing $n$ also divides $a$ then 
  \[
    \left|\disc\left( \Q\left( \alpha \right) \right)\right| \leq
    \prod_{p\mid a} p^{n v_p\left( a \right)} =
    \left| a^n \right| \leq 22^n
  \]
  However, according to \cite{OdlyzkoBoundsII} for any large enough $n$
  \[
    \left|\disc\left( \Q\left( \alpha \right) \right) \right| > 22^n
  \]
  (or $> 44^n$ assuming GRH) and hence there is a prime dividing $\disc\left( \Q\left( \alpha \right) \right)$ that does not divide $a$.
  The consequence follows from Lemma \ref{lem:inertia_group_acts_as_transposition}.
\end{proof}

\begin{lemma}
  \label{lem:bound_on_disc_of_intermediate_extensions}
  Assume that $\left|a\right| > 2$ and that $n \in S\left( \left| a \right| \right)$.
  Then for any intermediate extension $\Q\left( \alpha \right) / K / \Q$ of degree $m$, 
  \[
    \left| \disc\left( K \right)^{1/m} \right| \leq \left| a \right|
  \]
\end{lemma} 
\begin{proof}
  Indeed, 
  \[
    \disc\left( \Q\left( \alpha \right) \right) = \disc\left( K \right)^{\left[ \Q\left( \alpha \right) : K \right]} \Norm_{K / \Q}\left( \disc\left( \Q\left( \alpha \right) / K \right) \right)
  \]
  and hence for any prime $p$ dividing $\disc\left( K \right)$, 
  \[
    \frac{n}{m} v_p\left( \disc\left( K \right) \right) \leq  
    v_p\left( \disc\left( \Q\left( \alpha \right) \right) \right) \leq
    v_p\left( \disc\left( f_{n,a} \right) \right) = 
    n \min\left\{ v_p\left( n \right), v_p\left( a \right) \right\} \leq
    n v_p\left( a \right)
  \]
  and hence 
  \[
    \frac{v_p\left( \disc\left( K \right) \right)}{m} \leq v_p\left( a \right)
  \]
  Since it is true for any prime dividing $\disc\left( K \right)$, it follows that
  \[
    \left| \disc\left( K \right)^{1/m} \right| \leq \left| a \right|
  \]
\end{proof}
\begin{lemma}
  \label{lem:if_no_number_field_with_small_root_disc_then_galois_group_is_full}
  Let $n \in S\left( \left| a \right| \right)$.
  Assume that there is no intermediate extension $\Q \subsetneq K \subsetneq \Q\left( \alpha \right)$ that is ramified only over primes dividing $a$ and whose root discriminant is at most $\left| a \right|$.
  Then $\Gal\left( f_{n,a} \right) = S_n$. 
\end{lemma}
\begin{proof}
  Assume that $n$ satisfies the above conditions.
  Then by Lemma \ref{lem:bound_on_disc_of_intermediate_extensions}, $\Q\left( \alpha \right)$ has no proper intermediate extensions.
  Therefore, $\Hom\left( \Q\left( \alpha \right), \overline{\Q} \right)$ has no $G$-equivariant quotients and hence $G$ acts primitively on $\Hom\left( \Q\left( \alpha \right), \overline{\Q} \right)$.
  However, by Lemma \ref{lem:inertia_group_acts_as_transposition} and Lemma \ref{lem:there_is_a_prime_dividing_disc_but_not_a}, there is a prime $p$ such that $I\left( p \right)$ acts on $\Hom\left( \Q\left( \alpha \right), \overline{\Q} \right)$ as a single transposition and hence $G$ contains a transposition.
  Therefore, since any primitive subgroup of $S_n$ that contains a transposition is $S_n$, it follows that $G = S_n$.
\end{proof}

\editorial{In the following I am assuming that the Odlyzko bound on root discriminants is strictly increasing, even though I could not find a reference for a prood of this fact (and Gemini claims there isn't one).}

We can now use calculated minimal root discriminants of number fields ramified only over primes dividing $a$ to deduce which values of $a$ satisfy the condition of lemma \ref{lem:if_no_number_field_with_small_root_disc_then_galois_group_is_full}. 

\editorial{I used \url{https://www.lmfdb.org/NumberField/} to compose the following corollaries, I was not sure how to cite it.}

\begin{corollary}
  For any $n \in S\left( 3 \right)$, $\Gal\left( f_{n,\pm3} \right) = S_n$.  
\end{corollary}
\begin{proof}
  According to \url{https://www.lmfdb.org/NumberField/?ram_quantifier=subset&ram_primes=3&rd=-3-3}, the only non-trvial number field that is ramified only over $3$ and whose root discriminant is at most $3$ is the cyclotomic extension $\Q\left( \zeta_3 \right)$.
\end{proof}


\end{document}